\bookchapter{Bab 52: Empat Komponen dan Kitaran Hayat Activity}

\chapterintro{Jika Compose ialah 'wajah' aplikasi, bab ini pula bercakap tentang 'rangka' aplikasi. Sistem Android ada empat komponen asas; semua aplikasi dibina daripadanya; dan yang paling penting, Activity, mempunyai satu 'rajah kitaran hayat' yang wajib terpahat dalam kepala setiap pembangun Android. Selepas bab ini, anda akan memahami teka-teki abadi seperti 'kenapa bila masuk latar belakang dan kembali semula, keadaan aplikasi saya hilang'.}

\booksection{Empat Komponen: 'Tiang Seri' Android}

Android secara rasmi mentakrifkan empat komponen asas, masing-masing bertanggungjawab atas satu jenis kerja:

\begin{center}
\begin{tabularx}{\linewidth}{|X|X|X|}
\hline
\textbf{Komponen} & \textbf{Tanggungjawab} & \textbf{Ibaratnya} \tabularnewline
\hline
Activity & Satu antara muka yang boleh dilihat & Pemuka kedai — boleh dilihat dan boleh berinteraksi oleh pengguna \tabularnewline
Service & 'Pekerja halimunan' yang bekerja di latar belakang & Pekerja gudang — memainkan muzik, memuat turun fail, tidak perlukan antara muka \tabularnewline
BroadcastReceiver & Menerima siaran sistem & Rumah pengawal — 'pembesar suara' yang menyiarkan 'bateri habis', 'rangkaian terputus' \tabularnewline
ContentProvider & Berkongsi data merentas aplikasi & Kabinet fail awam — contohnya berkongsi kenalan dengan aplikasi lain \tabularnewline
\hline
\end{tabularx}
\end{center}

\begin{center}\small\color{gray}Sekilas pandang empat komponen Android\end{center}

Pada peringkat pemula, 95\% masa anda dihabiskan berurusan dengan Activity; Service dan BroadcastReceiver digunakan semasa memainkan muzik atau menghantar pemberitahuan; ContentProvider paling jarang digunakan (sistem sudah membungkus data yang lazim). Juga perhatikan: keempat-empat komponen mesti 'mendaftar' dalam AndroidManifest.xml; tanpa pendaftaran, sistem tidak jumpa mereka.

\booksection{Kitaran Hayat Activity: 'Peta Harta Karun' Pertama Pembangunan Android}

Aplikasi pada telefon boleh terganggu bila-bila masa: panggilan telefon masuk, pengguna bertukar ke WeChat untuk membalas mesej, memori sistem tidak cukup lalu mematikan anda… Sepanjang hayat Activity, ia melawati rajah ini secara berulang:

onCreate (lahir: pemulaan)

onStart (sudah kelihatan)

onResume (aktif di hadapan)

onPause (sebahagian tertutup)

onStop (langsung tidak kelihatan)

onDestroy (dimusnahkan)

onRestart / onStart / onResume: kembali ke hadapan melalui jalan yang sama

\begin{center}\small\color{gray}Rajah 4: Kitaran hayat Activity (arah ke hadapan: baris atas lahir baris bawah keluar; semasa kembali, berbalik melalui jalan asal)\end{center}

\begin{center}
\begin{tabularx}{\linewidth}{|X|X|X|}
\hline
\textbf{Panggil balik} & \textbf{Bila dicetuskan} & \textbf{Kegunaan biasa} \tabularnewline
\hline
onCreate & Semasa Activity dicipta (sekali) & Memuatkan antara muka, memulakan data \tabularnewline
onStart & Antara muka menjadi kelihatan & Memulakan beberapa logik ringan \tabularnewline
onResume & Mendapat fokus, boleh berinteraksi & Memulakan animasi, menyambung semula main balik \tabularnewline
onPause & Dilindungi sebahagian / sejurus bertukar pergi & Jeda animasi, simpan draf \tabularnewline
onStop & Sama sekali tidak kelihatan & Melepaskan sumber berat, menjeda muat turun \tabularnewline
onDestroy & Sebelum dimusnahkan & Pembersihan terakhir \tabularnewline
\hline
\end{tabularx}
\end{center}

\begin{center}\small\color{gray}Kegunaan biasa enam panggil balik kitaran hayat\end{center}

\textbf{Amaran Perangkap: Memusingkan Skrin Akan 'Membunuh' Activity}\par

Adegan terjebak paling klasik pemula: aplikasi dipusingkan sedikit, semua teks dalam kotak input hilang. Punca: memusingkan skrin secara lalai akan memusnahkan dan mencipta semula Activity (melalui sekali onDestroy onCreate). Penyelesaian: guna rememberSaveable untuk menyimpan keadaan (pengganti remember dalam Compose); ia akan mengembalikan nilai secara automatik semasa dibina semula. Akan digunakan berulang kali dalam Bab 6 dan 7.

\booksection{Intent: 'Tiket Perjalanan' Masuk Keluar Aplikasi}

Intent (niat) ialah 'tiket' komunikasi antara komponen. Mahu melompat dari satu halaman ke halaman lain, atau membuka pelayar, atau membangkitkan kamera, semuanya perlu 'beli tiket dan naik':

\begin{lstlisting}[language=Kotlin]
// Cara satu: Intent eksplisit—nyatakan halaman destinasi dengan jelas
val intent = Intent(this, DetailActivity::class.java)
intent.putExtra("name", "Zhang San")     // Seperti menitipkan "bagasi" dalam kenderaan (data)
startActivity(intent)

// Cara dua: Intent implisit—cukup sebut mahu buat apa, sistem carikan pelakunya
val web = Intent(Intent.ACTION_VIEW, Uri.parse("https://developer.android.com"))
startActivity(web)   // Sistem akan buka pelayar/atau minta anda pilih
\end{lstlisting}

Kedua-dua cara ada kegunaan masing-masing: eksplisit untuk lompatan halaman di dalam aplikasi sendiri; implisit untuk memanggil kebolehan sistem atau aplikasi lain (membuka laman web, membuat panggilan, berkongsi). Dalam projek Compose, dalam Composable gunakan LocalContext.current untuk mendapatkan Context maka boleh startActivity.

\booksection{Berbilang Halaman dan Navigasi}

Aplikasi sebenar semuanya ada berbilang halaman. Penyelesaian navigasi yang disyorkan secara rasmi untuk Compose ialah Navigation Compose (pustaka kebergantungan androidx.navigation):

\begin{lstlisting}[language=Kotlin]
val navController = rememberNavController()

NavHost(navController, startDestination = "home") {
    composable("home") { HomeScreen(navController) }
    composable("detail/{name}") { entry ->
        val name = entry.arguments?.getString("name")
        DetailScreen(name ?: "Tidak diketahui")
    }
}

// Berpindah halaman dalam halaman:
navController.navigate("detail/Zhang San")     // dengan parameter
navController.popBackStack()              // kembali ke halaman sebelumnya
\end{lstlisting}

\begin{tcbinfo}{Petua: Kalau Halaman Makin Banyak Macam Mana?}
Apabila projek membesar, kumpulkan semua laluan halaman dalam satu fail (contohnya AppNavGraph.kt) untuk diurus secara pusat; jangan biarkan logik 'pergi ke mana' berselerak merentas halaman. Ini langkah pertama dari 'boleh berjalan' ke 'boleh diselenggara'.

\end{tcbinfo}

\booksection{Sambil Membiasa Diri: Apa Itu Fragment (dan Kenapa Buku Ini Tidak Mengajarnya)}

Semasa merujuk tutorial lama, anda pasti akan bertemu Fragment — ia 'cekak antara muka boleh guna semula' dalam Activity, wira utama dalam seni bina berbilang halaman lama. Tetapi di era Compose, halaman cukup menggunakan fungsi Composable + laluan Navigation, dan pihak rasmi pun sudah melemahkan kedudukan Fragment. Maka buku ini tidak memasukkannya ke dalam laluan utama.

Cadangan untuk anda: projek baharu jangan aktif menggunakan Fragment; tetapi semasa membaca kod lama atau apabila ditanya dalam temu duga, tahu bahawa ia 'cekak antara muka yang disarang dalam Activity' sudah memadai. Sumber pembelajaran sentiasa utamakan yang selepas 2023 — 'kalendar lama' tindanan teknologi Android berubah lebih pantas daripada disangka.

\begin{tcbpractice}{Mari cuba: Beri Pengira Itu 'Ingatan'}
Ubah suai pengira Bab 51: (1) tukar remember kepada rememberSaveable, pusingkan skrin dan sahkan nombor masih ada; (2) tambah butang 'Tentang', guna Navigation untuk pergi ke halaman kedua yang memaparkan nombor versi; (3) cetak log (Log.d) dalam onCreate/onPause kedua-dua halaman, dan perhatikan susunan panggilan kitaran hayat dengan Logcat. Selepas tiga langkah ini, kitaran hayat bukan lagi teori di atas kertas.

\end{tcbpractice}

\begin{tcbsummary}{Ringkasan Bab}
\begin{itemize}
  \item Empat komponen: Activity (pemuka), Service (latar belakang), BroadcastReceiver (pembesar suara), ContentProvider (kabinet fail)
  \item Enam panggil balik kitaran hayat: onCreate onStart onResume onPause onStop onDestroy
  \item Memusingkan skrin memusnahkan dan membina semula Activity: guna rememberSaveable untuk menyelamatkan keadaan
  \item Intent ialah 'tiket': eksplisit melompat halaman sendiri, implisit memanggil kebolehan sistem
  \item Berbilang halaman guna Navigation Compose: laluan NavHost + composable
\end{itemize}

\end{tcbsummary}
